% TOP_PROBLEM: {"id": 30006866, "title": "Sinai's Positive Metric Entropy Conjecture for the Chirikov Standard Map", "category_id": 11, "category": {"id": 11, "name": "dynamical_systems", "display_name": "Dynamical Systems", "description": "Problems about long-term behavior of deterministic systems, Hamiltonian dynamics, and periodic orbits.", "slug": "dynamical-systems", "order_index": 11, "created_at": "2026-07-31T15:26:25.671Z"}, "status": "open"}
% FIRST_POSTED: 2026-09-27
% !TeX program = lualatex
% ATTEMPT_STATUS: unresolved
% Model: GPT-6 Astra Ultra; continuation dated 27 September 2026.
\documentclass[11pt]{article}
\usepackage[margin=25mm]{geometry}
\usepackage{fontspec}
\setmainfont{Latin Modern Roman}
\usepackage{amsmath,amssymb,mathtools}
\usepackage{unicode-math}
\setmathfont{Latin Modern Math}
\usepackage{xurl,hyperref}
\hypersetup{colorlinks=true,urlcolor=blue}
\setcounter{secnumdepth}{0}
\setlength{\parindent}{0pt}
\setlength{\parskip}{5pt}
\setlength{\emergencystretch}{3em}
\begin{document}
\section*{\texorpdfstring{Sinai's Positive Metric Entropy Conjecture for the Chirikov Standard Map}{Sinai's Positive Metric Entropy Conjecture for the Chirikov Standard Map}}\label{30006866}
\subsection{Definitions and mathematical statement}
On $\mathbb T^2={\mathbb{R}}^2/{\mathbb{Z}}^2$, define the area-preserving standard map
\[
 f_k(x,y)=\bigl(x+y+k\sin(2\pi x),\ y+k\sin(2\pi x)\bigr)\pmod{{\mathbb{Z}}^2}.
\]
Let $m$ be normalized Lebesgue area. Its metric entropy $h_m(f_k)$ measures the asymptotic information rate of measurable partitions under iteration, taking the supremum over finite partitions. The conjectured assertion is
\[
 \operatorname{Leb}\{k\in{\mathbb{R}}:h_m(f_k)>0\}>0.
\]
For this smooth area-preserving map, the entropy formula identifies positivity with a positive-$m$-measure set of points having positive maximal Lyapunov exponent, where that exponent is the long-time exponential growth rate of tangent vectors. Positive topological entropy on a zero-area invariant set is not sufficient. The parameter normalization here is exactly the one in the catalogue.
\par\smallskip\textit{Formulation and concept references:} \href{https://www.math.uci.edu/~asgor/Psim076.pdf}{[S1]}.

\subsection{Short English statement}
For a positive-measure set of parameters, does the standard map exhibit chaotic tangent growth on a positive-area set of trajectories?
\subsection{Research attempt}
\textbf{Target and source audit.}
The measure in this question is the fixed area measure $m$, not a measure
supported on a horseshoe.  Both the phase-space and the parameter-space
positivity requirements must be met.  The primary paper [E1] constructs
positive metric entropy by conservative perturbations, including perturbations
near standard maps.  Its perturbations are not required to stay in the
one-parameter family $f_k$.  The more recent standard-map result [E2] concerns
topological entropy.  Neither statement supplies the displayed target.  This
attempt gives explicit tangent estimates and identifies the precise infinite-time
measure estimate that those estimates fail to provide.

\textbf{Normalization, elementary dynamics, and a zero-entropy endpoint.}
Put $a(x)=2\pi k\cos(2\pi x)$.  Direct differentiation gives
\[
 Df_k(x,y)=\begin{pmatrix}1+a(x)&1\\a(x)&1\end{pmatrix},\qquad
 \det Df_k=1.
\]
The inverse is obtained from $x=X-Y$, then
$y=Y-k\sin(2\pi(X-Y))$.  Thus the formula indeed defines a smooth area-preserving
diffeomorphism on the torus.  The translation $C(x,y)=(x+1/2,y)$ satisfies
$C^{-1}f_kC=f_{-k}$, so entropy is an even function of the parameter.
For $k=0$,
\[
 f_0^N(x,y)=(x+Ny,y),\qquad Df_0^N=\begin{pmatrix}1&N\\0&1\end{pmatrix}.
\]
The derivative grows at most linearly.  Every Lyapunov exponent is zero,
and the smooth area entropy formula stated in the problem gives $h_m(f_0)=0$.
Consequently a proof cannot consist of an argument claiming positivity for
all parameters without an exception.

The points $(0,0)$ and $(1/2,0)$ are fixed.  Their derivative traces are
$2+2\pi k$ and $2-2\pi k$.  For $k>0$ the first is hyperbolic.  For
$0<k<2/\pi$ the second is elliptic at the linear level.  Interchanging their
roles handles negative $k$.  This proves existence of exponentially growing
directions on a periodic orbit for every nonzero parameter.  A finite orbit
has area zero, so even this exact exponential growth contributes nothing to
the integral of the positive exponent against $m$.

\textbf{A quantitative cone calculation.}
Use the integral linear torus coordinates $(q,r)=(x,x-y)$.  In these coordinates
the map is
\[
 F_k(q,r)=(2q-r+k\sin(2\pi q),q),\qquad
 DF_k=\begin{pmatrix}A(q)&-1\\1&0\end{pmatrix},\quad
 A(q)=2+2\pi k\cos(2\pi q).
\]
Fix $K>2$ and let $G=\{|A(q)|\geq K\}$.  On the cone
$\mathcal C=\{(u,v):|v|\leq|u|\}$, with the maximum norm,
\[
 |Au-v|\geq (K-1)|u|,
 \qquad |u|\leq\frac{|Au-v|}{K-1}.
\]
Thus $DF_k$ maps $\mathcal C$ into its interior and expands each vector in
that cone by at least $K-1$.  The sign of $A$ is irrelevant to this estimate.
If all the first $N$ iterates of a point lie in $G$, induction gives
\[
 \|DF_k^N(q,r)(1,0)\|_\infty\geq(K-1)^N.
\]
The coordinate change is fixed and invertible, so it does not alter Lyapunov
exponents.

This yields a clean sufficient condition.  If a measurable invariant set
$E\subset\bigcap_{j\in\mathbb Z}F_k^{-j}G$ has $m(E)>0$, then almost every
point of $E$ has maximal exponent at least $\log(K-1)$, by Oseledets' theorem
and the displayed expansion.  The entropy formula would give
\[
 h_m(f_k)\geq m(E)\log(K-1)>0.
\]
The external ergodic theorems here apply because the map and its inverse are
smooth on a compact manifold.  The missing assertion is existence of such a
positive-area invariant set, not the cone estimate.

\textbf{Why the critical strip cannot just be discarded.}
Write $b=2\pi|k|$.  When $b\geq2(K+2)$, the bad strip
$B=\{|2+b\cos(2\pi q)|<K\}$ has area
\[
 \varepsilon_k=\frac1\pi\left[
 \arccos\!\left(\frac{-K-2}{b}\right)
 -\arccos\!\left(\frac{K-2}{b}\right)\right]
 \leq\frac{4K}{\pi\sqrt3\,b}.
\]
Indeed both arguments belong to $[-1/2,1/2]$, where the absolute derivative
of $\arccos$ is at most $2/\sqrt3$, and their separation is $2K/b$.
The sign change for negative $k$ is absorbed by translation in $q$.
Since area is preserved, the union bound gives
\[
 m\left(\bigcap_{j=0}^{N-1}F_k^{-j}G\right)
 \geq 1-N\varepsilon_k.
\]
For a large fixed parameter this certifies exponential tangent growth over a
long finite window for most initial points.  It becomes vacuous when
$N\varepsilon_k\geq1$.  Taking $k$ large after fixing $N$ does not prove an
infinite-time assertion at any fixed $k$.  Interchanging those quantifiers is
exactly the error an apparent elementary proof would make.

Even a small frequency of visits to $B$ is insufficient without control of
their effect on directions.  An expanding direction can rotate into a
contracting direction.  For example, with
\[
 D=\begin{pmatrix}\lambda&0\\0&\lambda^{-1}\end{pmatrix},\qquad
 R=\begin{pmatrix}0&-1\\1&0\end{pmatrix},\quad\lambda>1,
\]
one has $DRDR=-I$.  Each factor $DR$ has operator norm $\lambda$, while their
product has norm one.  This is not a model orbit of the standard map and is
not a counterexample to the conjecture.  It is an exact countercheck to the
algebraic inference that multiplying large one-step norms gives a lower bound
on the norm of the product.

\textbf{The direction of the subadditivity bound.}
Let
\[
 a_N=\int_{\mathbb T^2}\log\|Df_k^N(z)\|\,dm(z)
\]
for the Euclidean operator norm.  The chain rule and invariance of $m$ imply
$a_{N+M}\leq a_N+a_M$.  The subadditive ergodic theorem therefore identifies
the integral of the maximal exponent with
\[
 \int\lambda_+(z)\,dm(z)=\inf_{N\geq1}\frac{a_N}{N}.
\]
Here determinant one makes the maximal exponent nonnegative.  A computed
positive value of $a_N/N$, even if rigorously enclosed for a large $N$, is
an upper bound on this infimum.  It is not a lower bound.  To prove positivity
by this route one needs a uniform positive lower bound for all sufficiently
large $N$, or a separate mechanism controlling cancellations at successive
scales.  Sampling finite-time exponents and observing stable positive numbers
does not supply that uniform estimate.

There is also a parameter issue after phase-space positivity is established.
A proof for one special $k$ would be meaningful progress, but would not by
itself prove positive Lebesgue measure of such parameters.  Metric entropy is
not automatically lower semicontinuous in the required form.  A positive
entropy perturbation in the space of all area-preserving diffeomorphisms can
leave the standard-map curve.  Its existence cannot be converted into a set
of parameters by projecting a neighborhood onto that curve.

\textbf{A possible route and the first unproved step.}
One could try to supplement the good-strip cone estimate with a return map
that groups a visit to $B$ together with its recovery interval.  A sufficient
construction would give a positive-area invariant domain, an induced map
with integrable return time, and a measurable cone field expanded by a factor
larger than one on average over those returns.  The recovery estimates must
track the direction, not just the orbit's distance from the critical curves.
The elementary bounds above supply neither a positive-area domain nor
integrability and positive mean expansion for this induced system.  They also
do not make those properties persist for a positive-measure parameter set.

\textbf{Verification and outcome.}
The saved diagnostic checks determinant and trace identities, the cancellation
matrix product, and sample cone inequalities exactly for integer data.  It
also evaluates the critical-strip area formula numerically at selected
parameters.  These finite checks test constants and signs, not asymptotic
entropy.  The result of this attempt is a precise sufficient invariant-set
criterion and an explanation of why predominantly expanding one-step
derivatives fail to verify it.  The required positive-measure set of parameters
has not been constructed.
\subsection{Sources}
\begin{itemize}
\item[S1]A. Gorodetski and V. Kaloshin, "Conservative Homoclinic Bifurcations and Some Applications," Proc. Steklov Inst. Math. 267 (2009), 76-90. DOI:10.1134/S0081543809040063.
\url{https://www.math.uci.edu/~asgor/Psim076.pdf}.
\textit{Formulation source.}
\item[E1]Pierre Berger and Dimitry Turaev, \emph{On Herman's Positive Entropy Conjecture}.
\url{https://arxiv.org/abs/1704.02473}.
\textit{Primary perturbation result; perturbations are not confined to the standard-map parameter family.}
\item[E2]Fernando Oliveira, \emph{Positive topological entropy for the Standard Map} (2024).
\url{https://arxiv.org/abs/2404.02209}.
\textit{Primary topological-entropy result; it does not establish positivity for Lebesgue metric entropy.}
\end{itemize}
\end{document}
